% \newpage
\subsection*{A2: Retrieval-Augmented Code Generation and Question Answering} \label{sec:app:retrieve}


Retrieval augmentation has emerged as a practical and effective approach for mitigating the intrinsic limitations of LLMs by incorporating external documents.  
In this section, we employ \libName to build a Retrieval-Augmented Generation (RAG) system~\cite{lewis2020retrieval, parvez2021retrieval} named Retrieval-augmented Chat. 
The system consists of two agents: a Retrieval-augmented User Proxy agent and a Retrieval-augmented Assistant agent, both of which are extended from built-in agents from \libName. The Retrieval-augmented User Proxy includes a vector database~\cite{chromadb} with SentenceTransformers~\cite{reimers-2019-sentence-bert} as the context retriever. A detailed workflow description of the Retrieval-augmented Chat is provided in Appendix~\ref{append:application}.

We evaluate Retrieval-augmented Chat in both question-answering and code-generation scenarios. (\textbf{Scenario 1}) We first perform an evaluation regarding natural question answering on the Natural Questions dataset~\cite{kwiatkowski2019natural} and report results in Figure~\ref{fig:res:retri}. In this evaluation, we compare our system with DPR (Dense Passage
Retrieval) following an existing evaluation\footnote{The results of DPR with GPT-3.5 shown in Figure~\ref{fig:res:retri} are from \cite{adlakha2023evaluating}. We use GPT-3.5 as a shorthand for GPT-3.5-turbo.} practice~\cite{adlakha2023evaluating}. Leveraging the conversational design and natural-language control,  \libName introduces a novel \emph{interactive retrieval} feature in this application: whenever the retrieved context does not contain the information, instead of terminating, the LLM-based assistant would reply ``\emph{Sorry, I cannot find any
information about... UPDATE CONTEXT.}” which will invoke more retrieval attempts. We conduct an ablation study in which we prompt the assistant agent to say \emph{``I don't know"} instead of \emph{``UPDATE CONTEXT."} in cases where relevant information is not found, and report results in Figure~\ref{fig:res:retri}. The results show that the interactive retrieval mechanism indeed plays a non-trivial role in the process. We give a concrete example and results using this appealing feature in Appendix~\ref{appendix_sec:app:retrive_chat}. (\textbf{Scenario 2}) We further demonstrate how Retrieval-augmented Chat aids in generating code based on a given codebase that contains code not included in GPT-4's training data. Evaluation and demonstration details for both scenarios are included in Appendix~\ref{appendix_sec:app:retrive_chat}.  